% ============================================================================
% LANGENTWURF LE-01c: Artikel + Alphabet
% Spanisch Klasse 7 - Sequenz 1: Empezamos
% ============================================================================
% Tier: 1 (vollständig)
% Doppelstunde: 1c (Std. 5-6 von 8)
% Fachfarbe: #c41e3a (Karminrot)
% ============================================================================

\documentclass[11pt,a4paper]{article}

\usepackage[utf8]{inputenc}
\usepackage[T1]{fontenc}
\usepackage[ngerman]{babel}
\usepackage{geometry}
\usepackage{fancyhdr}
\usepackage{lastpage}
\usepackage{xcolor}
\usepackage{tcolorbox}
\usepackage{enumitem}
\usepackage{tabularx}
\usepackage{longtable}
\usepackage{array}
\usepackage{booktabs}
\usepackage{amssymb}

% ============================================================================
% KONFIGURATION
% ============================================================================

\newcommand{\fach}{Spanisch}
\newcommand{\fachfarbe}{c41e3a}
\newcommand{\jahrgangsstufe}{7}
\newcommand{\schulart}{Gesamtschule}
\newcommand{\stundenthema}{Artikel + Alphabet}
\newcommand{\reihenthema}{Empezamos}
\newcommand{\stundennummer}{1c}
\newcommand{\reihenstunden}{4}
\newcommand{\gession}{A1}
\newcommand{\lehrwerk}{Qué pasa 1}
\newcommand{\lehrwerkseite}{S. 16-21}
\newcommand{\lerngruppengroesse}{24}

% ============================================================================
% LAYOUT
% ============================================================================
\geometry{left=2.5cm, right=2.5cm, top=2.5cm, bottom=2.5cm}
\definecolor{fach}{HTML}{\fachfarbe}
\definecolor{gray}{HTML}{666666}
\definecolor{lightgray}{HTML}{f5f5f5}
\definecolor{successgreen}{HTML}{2a7a4b}
\definecolor{warningorange}{HTML}{b35c00}
\definecolor{basis}{HTML}{2a7a4b}
\definecolor{standard}{HTML}{2c5aa0}
\definecolor{erweiterung}{HTML}{7b1fa2}

\tcbuselibrary{skins,breakable}

\newtcolorbox{infobox}[1][]{
    colback=lightgray,
    colframe=fach,
    fonttitle=\bfseries,
    title=#1,
    breakable
}

\newtcolorbox{scorebox}{
    colback=white,
    colframe=fach,
    fonttitle=\bfseries,
    title=Didaktik-Score,
    breakable
}

\pagestyle{fancy}
\fancyhf{}
\fancyhead[L]{\small\textcolor{gray}{\fach{} -- Jg. \jahrgangsstufe{} -- \stundenthema}}
\fancyhead[R]{\small\textcolor{gray}{LE-\stundennummer{} | Tier 1}}
\fancyfoot[C]{\small\thepage{} / \pageref{LastPage}}
\renewcommand{\headrulewidth}{0.4pt}

% Differenzierungsmarker
\newcommand{\B}{\textcolor{basis}{\textbf{[B]}}}
\newcommand{\Std}{\textcolor{standard}{\textbf{[S]}}}
\newcommand{\E}{\textcolor{erweiterung}{\textbf{[E]}}}

% ============================================================================
% DOKUMENT
% ============================================================================
\begin{document}

% ============================================================================
% DECKBLATT
% ============================================================================
\begin{titlepage}
    \centering
    \vspace*{2cm}
    {\large\textcolor{gray}{\schulart}}\\[2cm]
    {\Huge\textcolor{fach}{\textbf{Langentwurf}}}\\[0.5cm]
    {\large LE-\stundennummer{} | Tier 1 (vollständig)}\\[2cm]
    {\LARGE\textbf{\stundenthema}}\\[0.5cm]
    {\large Sequenz: \reihenthema}\\[2cm]
    \begin{tabular}{rl}
        \textbf{Fach:} & \fach{} (A1) \\[0.3cm]
        \textbf{Klasse:} & \jahrgangsstufe \\[0.3cm]
        \textbf{Lehrwerk:} & \lehrwerk, \lehrwerkseite \\[0.3cm]
        \textbf{Doppelstunde:} & \stundennummer{} (Std. 5-6 von 8) \\[0.3cm]
    \end{tabular}
    \vfill
\end{titlepage}

\tableofcontents
\newpage

% ============================================================================
% 1. BEDINGUNGSANALYSE
% ============================================================================
\section{Bedingungsanalyse}

\subsection{Lerngruppe}
\begin{tabular}{ll}
    Kursgröße: & \lerngruppengroesse{} SuS \\
    Jahrgangsstufe: & \jahrgangsstufe \\
    Sprachniveau: & A1 (GER) -- Anfänger \\
    Fremdsprache: & 2. Fremdsprache (nach Englisch) \\
\end{tabular}

\subsection{Vorwissen aus 01a/01b}
\begin{itemize}
    \item Begrüßungen und Verabschiedungen
    \item Sich vorstellen (Name, Herkunft)
    \item Verb \textit{ser} (yo soy, tú eres, él/ella es, nosotros somos)
    \item Fragewörter: ¿Cómo?, ¿De dónde?
    \item Redemittel: Me llamo..., Soy de...
\end{itemize}

\subsection{Differenzierung (intern)}
\begin{tabular}{|l|p{3.5cm}|p{3.5cm}|p{3.5cm}|}
\hline
& \B{} \textbf{Basis} & \Std{} \textbf{Standard} & \E{} \textbf{Erweiterung} \\
\hline
\textbf{Ziel} & BR: Rezeption & MR: Produktion mit Hilfe & GY: Freie Produktion \\
\hline
\textbf{Artikel} & Zuordnung mit Tabelle & Einsetzen selbstständig & + Regeln formulieren \\
\hline
\textbf{Alphabet} & Nachsprechen & Buchstabieren mit Hilfe & Frei buchstabieren \\
\hline
\end{tabular}

% ============================================================================
% 2. SACHANALYSE
% ============================================================================
\section{Sachanalyse}

\subsection{Sprachliche Strukturen}

\textbf{Bestimmte Artikel:}
\begin{center}
\begin{tabular}{|l|c|c|}
\hline
& \textbf{Singular} & \textbf{Plural} \\
\hline
\textbf{maskulin} & el (el chico) & los (los chicos) \\
\hline
\textbf{feminin} & la (la chica) & las (las chicas) \\
\hline
\end{tabular}
\end{center}

\textbf{Genus-Regeln (Tendenz):}
\begin{itemize}
    \item Maskulin: Nomen auf \textbf{-o} (el libro, el cuaderno)
    \item Feminin: Nomen auf \textbf{-a} (la mesa, la silla)
    \item Ausnahmen: el día, el mapa, la mano, la foto
\end{itemize}

\textbf{Spanisches Alphabet:}
\begin{itemize}
    \item 27 Buchstaben (+ ñ)
    \item Besondere Aussprache: c, g, h, j, ll, ñ, qu, rr, v, z
    \item Akzentregeln (rezeptiv): á, é, í, ó, ú
\end{itemize}

\subsection{Kontrastive Betrachtung}
\begin{itemize}
    \item \textbf{Positiv:} Artikel-Konzept aus Deutsch bekannt
    \item \textbf{Interferenz:} Deutsches Neutrum existiert nicht
    \item \textbf{Interferenz:} Genusfehler (el problema, nicht *la problema)
    \item \textbf{Alphabet:} H wird nicht gesprochen, J wie deutsches CH
\end{itemize}

% ============================================================================
% 3. DIDAKTISCHE ANALYSE
% ============================================================================
\section{Didaktische Analyse}

\subsection{Stellung in der Sequenz}
Dies ist die \textbf{3. Doppelstunde} (Std. 5-6 von 8) der Sequenz ``\reihenthema''.

\subsection{Kompetenzziele (Rahmenplan MV)}
\begin{itemize}
    \item \textbf{Sprachlich:} Genus und Numerus von Nomen erkennen und Artikel zuordnen
    \item \textbf{Kommunikativ:} Namen buchstabieren können (Sprechen)
    \item \textbf{Methodisch:} Wörterbucharbeit vorbereiten (Genus-Markierung verstehen)
\end{itemize}

\subsection{Legitimation}
\textbf{Rahmenplan Spanisch MV:}
\begin{itemize}
    \item Verfügen über sprachliche Mittel: Grammatik -- Artikel, Genus
    \item Sprechen: Wörter buchstabieren, Namen nennen
\end{itemize}

% ============================================================================
% 4. STUNDENZIELE
% ============================================================================
\section{Stundenziele}

\subsection{Grobziel}
Die SuS erweitern ihre \textbf{grammatische Kompetenz}, indem sie die bestimmten Artikel (el, la, los, las) Nomen zuordnen und \textbf{Namen buchstabieren} können.

\subsection{Feinziele (operationalisiert)}
\begin{tabular}{|c|p{8cm}|c|c|}
\hline
\textbf{Nr.} & \textbf{Die SuS können...} & \textbf{AFB} & \textbf{Diff.} \\
\hline
FZ1 & die vier bestimmten Artikel (el, la, los, las) benennen und Nomen nach Genus/Numerus zuordnen & I & alle \\
\hline
FZ2 & das spanische Alphabet korrekt aussprechen und Unterschiede zum Deutschen erkennen & I & alle \\
\hline
FZ3 & einfache Nomen mit dem richtigen Artikel verwenden und eigene Namen buchstabieren & II & alle \\
\hline
FZ4 & Genus-Regeln (-o/-a) anwenden und Ausnahmen identifizieren & II/III & \Std{}/\E{} \\
\hline
\end{tabular}

% ============================================================================
% 5. METHODISCHE ANALYSE
% ============================================================================
\section{Methodische Analyse}

\subsection{Medien}
\begin{itemize}
    \item Tafel/Whiteboard
    \item Lehrwerk: \lehrwerk, \lehrwerkseite
    \item Arbeitsblatt: AB-01c.pdf
    \item Audio: Track 8-9 (Alphabet-Song)
    \item Bildkarten: Nomen mit Artikel
    \item Optional: LP-01c.html (digital)
\end{itemize}

\subsection{Sozialformen}
\begin{itemize}
    \item Plenum: Einführung Artikel, Alphabet-Drill
    \item Partnerarbeit: Buchstabier-Spiel
    \item Einzelarbeit: AB-Aufgaben
\end{itemize}

% ============================================================================
% 6. VERLAUFSPLANUNG
% ============================================================================
\section{Verlaufsplanung}

\textbf{1. Stunde (45 Min.) -- Schwerpunkt: Artikel}

\begin{longtable}{|p{1cm}|p{1.5cm}|p{4cm}|p{3cm}|p{2cm}|p{2cm}|}
\hline
\textbf{Zeit} & \textbf{Phase} & \textbf{Lehrerhandeln} & \textbf{Schülerhandeln} & \textbf{SF} & \textbf{Medien} \\
\hline
\endhead

0--5 & Einstieg & Gegenstände zeigen: ``¿Qué es esto?'' & Vermutungen äußern & Plenum & Objekte \\
\hline

5--15 & Input & Artikel einführen (el/la), Tafelanschrieb: Tabelle & Mitschreiben, Nachsprechen & Plenum & Tafel \\
\hline

15--25 & Übung 1 & Bildkarten hochhalten, Artikel abfragen & Chorisch antworten: ``el libro'', ``la mesa'' & Plenum & Karten \\
\hline

25--35 & Übung 2 & AB-01c Aufg. 1 erklären & Zuordnung: Nomen $\rightarrow$ Artikel & EA & AB \\
\hline

35--40 & Sicherung & Lösungen besprechen, Regeln (-o/-a) ableiten & Vergleichen, Regeln formulieren & Plenum & Tafel \\
\hline

40--45 & Über\-leitung & Plural einführen (los/las), Ausblick & Notieren & Plenum & Tafel \\
\hline
\end{longtable}

\vspace{1em}

\textbf{2. Stunde (45 Min.) -- Schwerpunkt: Alphabet}

\begin{longtable}{|p{1cm}|p{1.5cm}|p{4cm}|p{3cm}|p{2cm}|p{2cm}|}
\hline
\textbf{Zeit} & \textbf{Phase} & \textbf{Lehrerhandeln} & \textbf{Schülerhandeln} & \textbf{SF} & \textbf{Medien} \\
\hline
\endhead

0--5 & Aktivierung & Artikel-Quiz: ``el oder la?'' & Schnell antworten & Plenum & -- \\
\hline

5--15 & Input & Alphabet vorsprechen, Audio abspielen & Mitsprechen, auf Unterschiede achten & Plenum & Audio \\
\hline

15--25 & Übung 1 & Buchstabier-Drill: Wörter buchstabieren & Nachsprechen: ``hache-o-ele-a'' & Plenum & Tafel \\
\hline

25--35 & Übung 2 & PA: Namen buchstabieren ``¿Cómo se escribe...?'' & Partner fragt, anderer buchstabiert & PA & -- \\
\hline

35--40 & Sicherung & Merkkasten AB-01c gemeinsam ausfüllen & Übertragen & EA & AB \\
\hline

40--45 & Abschluss & Zusammenfassung, HA erklären & Aufschreiben & Plenum & -- \\
\hline
\end{longtable}

\textbf{Legende:} EA = Einzelarbeit, PA = Partnerarbeit, SF = Sozialform

% ============================================================================
% 7. ERWARTETE SCHWIERIGKEITEN
% ============================================================================
\section{Erwartete Schwierigkeiten}

\begin{tabular}{|p{5cm}|p{8cm}|}
\hline
\textbf{Schwierigkeit} & \textbf{Reaktion} \\
\hline
Genusfehler (*la libro) & Kontrastive Übung: el $\rightarrow$ -o, la $\rightarrow$ -a; Ausnahmen explizit thematisieren \\
\hline
Aussprache H (stumm) & Bewusstmachung: H = stumm im Spanischen; Drill: hola, hora, hombre \\
\hline
Verwechslung J/G & Minimalpaar-Übung: gente vs. jefe; G vor e/i = CH-Laut \\
\hline
Buchstabieren zu langsam & Mehr Übungszeit, Alphabet-Song wiederholen \\
\hline
Plural-Bildung unklar & Einfache Regel: + s (nach Vokal), + es (nach Konsonant) \\
\hline
\end{tabular}

% ============================================================================
% 8. HAUSAUFGABE
% ============================================================================
\section{Hausaufgabe}

\begin{itemize}
    \item Lehrwerk S. 18, Übung 4 (Artikel einsetzen)
    \item Vokabeln lernen: Artikel + 10 Nomen
    \item Optional: LP-01c.html bearbeiten (Buchstabier-Übung)
\end{itemize}

% ============================================================================
% 9. LÖSUNGEN
% ============================================================================
\section{Lösungen AB-01c}

\textbf{Merkkasten:}
\begin{itemize}
    \item maskulin Singular: el
    \item feminin Singular: la
    \item maskulin Plural: los
    \item feminin Plural: las
\end{itemize}

\textbf{Aufgabe 1 (Zuordnung):} (4P)
\begin{itemize}
    \item el: libro, cuaderno, bolígrafo, lápiz
    \item la: mesa, silla, pizarra, ventana
\end{itemize}

\textbf{Aufgabe 2 (Einsetzen):} (4P)
\begin{enumerate}[label=\alph*)]
    \item \textbf{el} chico
    \item \textbf{la} profesora
    \item \textbf{los} libros
    \item \textbf{las} chicas
\end{enumerate}

\textbf{Aufgabe 3 (Buchstabieren):} (6P)
\begin{itemize}
    \item HOLA: hache - o - ele - a
    \item GRACIAS: ge - erre - a - ce - i - a - ese
    \item [eigener Name]: individuell
\end{itemize}

\textbf{Aufgabe 4 (Dialog):} (6P)
\begin{itemize}
    \item A: ¿Cómo te llamas?
    \item B: Me llamo [Name].
    \item A: ¿Cómo se escribe?
    \item B: [Buchstabiert den Namen]
\end{itemize}

% ============================================================================
% 10. DIDAKTIK-SCORE
% ============================================================================
\newpage
\section{Didaktik-Score}

\begin{scorebox}
\textbf{Bewertung der didaktischen Qualität nach 10 Dimensionen}\\
\textbf{Ziel-Score: $\geq$ 9.0 (Premium-Qualität)}

\vspace{1em}
\begin{tabular}{|c|l|c|c|p{4.5cm}|}
\hline
\textbf{\#} & \textbf{Dimension} & \textbf{Gew.} & \textbf{Score} & \textbf{Notiz} \\
\hline
1 & Differenzierung & 15\% & 9/10 & [B] Zuordnung mit Hilfe, [S] selbstständig, [E] Regeln formulieren \\
\hline
2 & Taxonomie (AFB) & 10\% & 9/10 & AFB I: 40\%, AFB II: 50\%, AFB III: 10\% \\
\hline
3 & Lernziel-Alignment & 10\% & 10/10 & Rahmenplan: Artikel, Buchstabieren \\
\hline
4 & Aktivierung & 10\% & 9/10 & Chorisch, PA-Buchstabieren, Spiel \\
\hline
5 & Scaffolding & 10\% & 9/10 & Tabelle $\rightarrow$ Übung $\rightarrow$ freie Anwendung \\
\hline
6 & Feedback-Qualität & 10\% & 9/10 & Elaboriertes Feedback in LP \\
\hline
7 & Sprachliche Klarheit & 10\% & 10/10 & Altersgerecht, kurze Sätze \\
\hline
8 & Motivation \& Relevanz & 10\% & 9/10 & Namen buchstabieren = Alltagsrelevanz \\
\hline
9 & Inklusion (UDL) & 10\% & 9/10 & Audio + visuell + kinästhetisch \\
\hline
10 & Praktikabilität & 5\% & 9/10 & 2x45 Min realistisch \\
\hline
\multicolumn{3}{|r|}{\textbf{GESAMT:}} & \textbf{9.2/10} & \textcolor{successgreen}{\textbf{FREIGABE}} \\
\hline
\end{tabular}
\end{scorebox}

\end{document}
